\documentclass[11pt]{article}
\usepackage[margin=1in]{geometry}
\usepackage{booktabs}
\usepackage{array}
\usepackage{tabularx}
\usepackage{siunitx}
\usepackage{hyperref}

\sisetup{round-mode=places,round-precision=2}

\title{Memory-Efficient Torch Inference for AlphaGenome}
\author{}
\date{July 3, 2026}

\begin{document}
\maketitle

\begin{abstract}
We evaluated a sequence of memory-saving inference strategies for a Torch
implementation of AlphaGenome using a recreated merged human brain development
ATAC checkpoint. The strategies combine bfloat16 execution, precomputed
standardized convolutions, Triton int8 weight-only one-dimensional convolutions,
an encoder path that omits unneeded intermediate tensors, a custom no-indices
Triton max-pooling kernel, and fused high-resolution encoder blocks. On
131 kb input windows at batch size 32, the most memory-efficient successful
configuration reduced observed peak device memory from 92.6 GiB to 25.2 GiB
relative to default Torch inference, while preserving differential Pearson to
within 0.001. On 1 Mb input windows at batch size 2, the default path exceeded
48 GiB, whereas the memory-efficient configurations were between 24.1 GiB and
33.4 GiB observed peak memory.
\end{abstract}

\section{Experimental setup}

All experiments used the recreated merged Torch checkpoint from the best
LoRA+LoCon finetuning run. We evaluated the human brain development ATAC head at
resolution 128 on full valid and test splits. For 131 kb input windows, batch
size was 32 and the evaluation contained 2162 examples. For 1 Mb input windows,
a valid/test-only target cache was built, batch size was 2, and the evaluation
contained 269 examples. Peak device memory was sampled with \texttt{nvidia-smi};
PyTorch allocated and reserved memory were recorded from CUDA allocator
statistics.

\section{Memory-saving strategies}

The evaluated strategies are cumulative unless otherwise indicated.
\begin{enumerate}
    \item \textbf{Bfloat16 parameters and activations.} The
    \texttt{bf16\_params} baseline applies the aggressive bfloat16 Torch dtype
    policy for inference.
    \item \textbf{Precomputed standardized convolutions and Triton int8
    weight-only Conv1d.} Standardized convolution weights are materialized once
    for inference and replaced by direct effective convolutions. Eligible wide
    convolutions are then replaced by Triton int8 weight-only Conv1d kernels with
    floating-point output activations.
    \item \textbf{No-intermediates encoder path.} When only 128 bp predictions
    are requested, U-Net skip tensors that are not consumed are not retained.
    \item \textbf{Triton no-indices max-pooling.} The encoder max-pooling layer
    is replaced by a Triton implementation for kernel size 2 and stride 2 that
    does not materialize PyTorch max-pool indices.
    \item \textbf{Fused DNA embedder block.} The full-resolution
    \texttt{encoder.dna\_embedder.block} is fused into one Triton operation
    implementing RMSBatchNorm, the JAX GELU approximation, and int8 weight-only
    Conv1d.
    \item \textbf{Fused first downsampling block.} The
    \texttt{encoder.down\_blocks.0} module is replaced by a wrapper using fused
    ConvBlocks and in-place residual accumulation in inference.
\end{enumerate}

\section{Results}

\subsection{131 kb input windows}

\begin{table}[h]
\centering
\caption{Full valid/test evaluation for 131 kb windows at batch size 32.}
\label{tab:vram-131kb}
\begin{tabular}{lrrrrrr}
\toprule
Strategy & Ex./s & Peak MiB & Alloc MiB & Reserved MiB & Valid $r$ & Test $r$ \\
\midrule
Default & 9.41 & 92631 & 63140 & 91928 & 0.79893 & 0.80891 \\
Bfloat16 & 12.02 & 66935 & 42415 & 66248 & 0.79839 & 0.80874 \\
Triton Conv1d & 8.49 & 43611 & 27879 & 42924 & 0.79842 & 0.80877 \\
No intermediates & 8.51 & 43611 & 25319 & 42924 & 0.79842 & 0.80877 \\
Triton pool & 8.78 & 34907 & 25319 & 34220 & 0.79842 & 0.80877 \\
Fused embedder & 8.03 & 32347 & 21735 & 31660 & 0.79845 & 0.80872 \\
Fused down0 & 8.11 & 31323 & 25319 & 30636 & 0.79829 & 0.80860 \\
Fused embedder+down0 & 7.45 & 25179 & 19174 & 24492 & 0.79834 & 0.80860 \\
\bottomrule
\end{tabular}
\end{table}

At 131 kb, the most memory-efficient configuration reduced observed peak memory
from 92.6 GiB to 25.2 GiB relative to default inference and from 66.9 GiB to
25.2 GiB relative to bfloat16 inference. The best throughput among configurations
below 48 GiB was the Triton-pool strategy, at 8.78 examples/s and 34.9 GiB.

\subsection{1 Mb input windows}

\begin{table}[h]
\centering
\caption{Full valid/test evaluation for 1 Mb windows at batch size 2.}
\label{tab:vram-1mb}
\begin{tabular}{lrrrrrr}
\toprule
Strategy & Ex./s & Peak MiB & Alloc MiB & Reserved MiB & Valid $r$ & Test $r$ \\
\midrule
Default & 0.721 & 52595 & 38507 & 51908 & 0.81825 & 0.82707 \\
Bfloat16 & 0.937 & 47597 & 25445 & 46910 & 0.81522 & 0.82640 \\
Triton Conv1d & 0.736 & 33359 & 22508 & 32672 & 0.81558 & 0.82662 \\
No intermediates & 0.749 & 29263 & 15452 & 28576 & 0.81558 & 0.82662 \\
Triton pool & 0.753 & 30799 & 15452 & 30112 & 0.81558 & 0.82662 \\
Fused embedder & 0.701 & 29519 & 15453 & 28832 & 0.81554 & 0.82658 \\
Fused down0 & 0.697 & 27215 & 15452 & 26528 & 0.81541 & 0.82643 \\
Fused embedder+down0 & 0.650 & 24143 & 15452 & 23456 & 0.81534 & 0.82634 \\
\bottomrule
\end{tabular}
\end{table}

At 1 Mb, default inference exceeded 48 GiB observed peak memory. Bfloat16
inference fit just under 48 GiB and was the fastest measured option. The
Triton/no-intermediates configurations reduced live allocated memory further,
to approximately 15.5 GiB, while remaining well below 48 GiB observed memory.
The fused kernels primarily reduced allocator-reserved and observed memory at
this longer context length, but also reduced throughput.

\section{Discussion}

The most robust practical configuration is the combination of bfloat16,
precomputed standardized convolutions, Triton int8 weight-only Conv1d, the
128 bp no-intermediates encoder path, and Triton no-indices max-pooling. This
configuration preserved differential Pearson relative to the bfloat16/Triton
Conv1d baseline and kept peak memory below 48 GiB for both tested input lengths.
The fused high-resolution blocks are useful when minimizing observed memory is
more important than maximizing throughput. For 1 Mb inputs, the tower and
attention computations dominate a larger share of runtime, so convolution
fusions provide less throughput benefit than at shorter context length.

\end{document}
